% One-loop diagrams of the collinear function J_1 for g -> q qbar, drawn with TikZ-Feynman.
% The vertices are placed by hand, so pdflatex suffices. Same set-up as figures/dps_diagram.tex.
% Crossed circles: the hard vertices H_1 and H_2 entered by the two daughters. Double lines: the Wilson
% lines of the two collinear fields. Compile with pdflatex; the Supplement includes the resulting PDF.
% The names are those of the entries in paper2/r0/theory/gqq_diagrams.py and paper2/r0/derivation.
% "mgluon" is the gluon coil of TikZ-Feynman reflected about its path. A line or panel that is the
% mirror image of another one under y -> -y uses it, so that the two are exact mirror images.
\documentclass[border=2pt]{standalone}
\pdfinfoomitdate=1 \pdftrailerid{} \pdfsuppressptexinfo=-1     % the same bytes at every build
\usepackage[compat=1.1.0]{tikz-feynman}
\tikzfeynmanset{warn luatex=false, arrow size=1.15pt,
  mgluon/.style={/tikz/draw=none, /tikz/decoration={name=none},
    /tikz/postaction={/tikz/draw, /tikz/decoration={coil, aspect=0.75, segment length=1.5mm}, /tikz/decorate=true}}}
\tikzset{wilson/.style={double, double distance=1.0pt, line width=0.45pt}}
% the two crossed circles, their Wilson lines, and the panel label
\newcommand{\operators}{%
  \vertex[crossed dot] (o1) at (3.0, 0.8) {};
  \vertex[crossed dot] (o2) at (3.0,-0.8) {};}
\newcommand{\wilsonlines}{%
  \draw[wilson] (o1) -- (3.93, 1.29);
  \draw[wilson] (o2) -- (3.93,-1.29);}
\newcommand{\panel}[1]{\node[anchor=north west] at (-0.5, 1.95) {#1};}
\begin{document}
{\small
\begin{tikzpicture}

% ================================================================ tree splitting and vertex corrections
% (a) tree splitting
\begin{feynman}
  \vertex (p) at (0.0, 0.0) {\(g(p_1)\)};
  \vertex (s) at (1.5, 0.0);
  \operators
  \diagram* {
    (p) -- [gluon] (s),
    (o2) -- [fermion] (s) -- [fermion] (o1),
  };
\end{feynman}
\wilsonlines
\node at (2.05, 0.75) {\(k\)};
\node at (1.95,-0.80) {\(p_1-k\)};
\node[anchor=west] at (3.15, 0.45) {\(H_1\)};
\node[anchor=west] at (3.15,-0.45) {\(H_2\)};
\panel{(a)}

% (b) Vqq: gluon exchanged between the quark and the antiquark
\begin{scope}[xshift=4.6cm]
\begin{feynman}
  \vertex (p) at (0.3, 0.0);
  \vertex (s) at (1.5, 0.0);
  \vertex (a) at (2.5, 0.533);
  \vertex (b) at (2.5, -0.533);
  \operators
  \diagram* {
    (p) -- [gluon] (s),
    (o2) -- [plain] (b) -- [fermion] (s) -- [fermion] (a) -- [plain] (o1),
    (a) -- [gluon] (b),
  };
\end{feynman}
\wilsonlines
\panel{(b)~Vqq}
\end{scope}

% (c) Vggq: three-gluon vertex, both gluons attached to the quark line
\begin{scope}[xshift=9.2cm]
\begin{feynman}
  \vertex (p) at (0.0, 0.0);
  \vertex (s) at (1.0, 0.0);
  \vertex (a) at (2.15, 0.60);
  \vertex (b) at (2.15, -0.60);
  \operators
  \diagram* {
    (p) -- [gluon] (s),
    (s) -- [gluon] (a),
    (s) -- [mgluon] (b),
    (o2) -- [fermion] (b) -- [fermion] (a) -- [fermion] (o1),
  };
\end{feynman}
\wilsonlines
\panel{(c)~Vggq}
\end{scope}

% ================================================================ Wilson line of H_1
% (d) W1q: Wilson line of H_1 to the quark
\begin{scope}[yshift=-3.8cm]
\begin{feynman}
  \vertex (p) at (0.3, 0.0);
  \vertex (s) at (1.5, 0.0);
  \vertex (a) at (2.29, 0.42);
  \vertex (w) at (3.71, 1.18);
  \operators
  \diagram* {
    (p) -- [gluon] (s),
    (o2) -- [fermion] (s) -- [fermion] (a) -- [plain] (o1),
    (w) -- [gluon, bend right=80] (a),
  };
\end{feynman}
\wilsonlines
\panel{(d)~W1q}
\end{scope}

% (e) W1qb: Wilson line of H_1 to the antiquark
\begin{scope}[xshift=4.6cm, yshift=-3.8cm]
\begin{feynman}
  \vertex (p) at (0.3, 0.0);
  \vertex (s) at (1.5, 0.0);
  \vertex (a) at (2.25, -0.40);
  \vertex (w) at (3.71, 1.18);
  \operators
  \diagram* {
    (p) -- [gluon] (s),
    (o2) -- [plain] (a) -- [fermion] (s) -- [fermion] (o1),
    (w) -- [gluon, bend left=25] (a),
  };
\end{feynman}
\wilsonlines
\panel{(e)~W1qb}
\end{scope}

% (f) W1g: Wilson line of H_1 to the parent gluon
\begin{scope}[xshift=9.2cm, yshift=-3.8cm]
\begin{feynman}
  \vertex (p) at (0.1, 0.0);
  \vertex (g) at (0.8, 0.0);
  \vertex (s) at (1.5, 0.0);
  \vertex (w) at (3.71, 1.18);
  \operators
  \diagram* {
    (p) -- [gluon] (g) -- [gluon] (s),
    (o2) -- [fermion] (s) -- [fermion] (o1),
    (w) -- [gluon, bend right=58] (g),
  };
\end{feynman}
\wilsonlines
\panel{(f)~W1g}
\end{scope}

% ================================================================ Wilson line of H_2 (mirror images of the second row)
% (g) W2qb: Wilson line of H_2 to the antiquark
\begin{scope}[yshift=-7.6cm]
\begin{feynman}
  \vertex (p) at (0.3, 0.0);
  \vertex (s) at (1.5, 0.0);
  \vertex (a) at (2.29, -0.42);
  \vertex (w) at (3.71, -1.18);
  \operators
  \diagram* {
    (p) -- [gluon] (s),
    (o2) -- [plain] (a) -- [fermion] (s) -- [fermion] (o1),
    (w) -- [mgluon, bend left=80] (a),
  };
\end{feynman}
\wilsonlines
\panel{(g)~W2qb}
\end{scope}

% (h) W2q: Wilson line of H_2 to the quark
\begin{scope}[xshift=4.6cm, yshift=-7.6cm]
\begin{feynman}
  \vertex (p) at (0.3, 0.0);
  \vertex (s) at (1.5, 0.0);
  \vertex (a) at (2.25, 0.40);
  \vertex (w) at (3.71, -1.18);
  \operators
  \diagram* {
    (p) -- [gluon] (s),
    (o2) -- [fermion] (s) -- [fermion] (a) -- [plain] (o1),
    (w) -- [mgluon, bend right=25] (a),
  };
\end{feynman}
\wilsonlines
\panel{(h)~W2q}
\end{scope}

% (i) W2g: Wilson line of H_2 to the parent gluon
\begin{scope}[xshift=9.2cm, yshift=-7.6cm]
\begin{feynman}
  \vertex (p) at (0.1, 0.0);
  \vertex (g) at (0.8, 0.0);
  \vertex (s) at (1.5, 0.0);
  \vertex (w) at (3.71, -1.18);
  \operators
  \diagram* {
    (p) -- [gluon] (g) -- [gluon] (s),
    (o2) -- [fermion] (s) -- [fermion] (o1),
    (w) -- [mgluon, bend left=58] (g),
  };
\end{feynman}
\wilsonlines
\panel{(i)~W2g}
\end{scope}

\end{tikzpicture}}
\end{document}
